% 第6章 银行公司业务与风控建模
% 智慧银行实验教程——AI驱动的金融科技实践

\chapter{第6章\quad 银行公司业务与风控建模}
\chapteroverview{
本章学习目标：
\begin{enumerate}
  \item 理解公司银行业务的授信流程与贷后管理
  \item 掌握企业信用风险评估的经典模型
  \item 学会构建贷款FTP定价模型
  \item 能设计贷后风险预警系统
  \item 了解供应链金融的AI应用
\end{enumerate}
}

% ============================================================
\section{公司银行业务理论}
% ============================================================

公司银行（Corporate Banking），又称对公业务，是商业银行为企业、政府及机构客户提供的金融服务总称。与零售银行服务个人客户不同，公司银行面对的是组织客户，其业务复杂度、单笔金额和风险特征与零售业务存在本质差异。

\subsection{公司银行业务全景}

公司银行的核心业务线包括：

\begin{table}[H]
\centering
\begin{tabular}{lp{3.5cm}p{5.5cm}}
\toprule
\textbf{业务线} & \textbf{代表产品} & \textbf{核心逻辑} \\
\midrule
对公存款 & 活期/定期/协定存款、通知存款 & 低成本企业资金，银行负债端基石 \\
对公贷款 & 流动资金贷款、项目贷款、银团贷款 & 核心资产业务，利差收入主要来源 \\
票据业务 & 银行承兑汇票、商业承兑汇票、贴现 & 供应链支付工具+短期融资 \\
贸易融资 & 信用证、保函、福费廷、保理 & 服务进出口贸易，风险缓释工具 \\
现金管理 & 资金池、代发薪、集团结算 & 绑定企业结算账户，沉淀低成本存款 \\
投资银行 & 债券承销、并购顾问、资产证券化 & 中间业务收入，高附加值服务 \\
\bottomrule
\end{tabular}
\end{table}

公司银行业务的核心特征是"金额大、频次低、关系深"：单笔贷款动辄数百万至数十亿元，但客户数量远少于零售客户；银企关系通常持续多年，客户转换成本高。

\subsection{授信全流程}

银行对公授信遵循严格的全流程管理，确保风险可控：

\begin{table}[H]
\centering
\begin{tabular}{lp{3cm}p{6cm}}
\toprule
\textbf{阶段} & \textbf{责任方} & \textbf{核心工作} \\
\midrule
贷前调查 & 客户经理 & 实地走访、收集资料、初步评估 \\
评审审查 & 风险管理部 & 独立审查、风险量化、方案优化 \\
审批决策 & 有权审批人 & 集体审议、一票否决、条件设定 \\
放款执行 & 放款中心 & 合同审核、抵质押登记、条件落实 \\
贷后管理 & 客户经理+风控 & 定期检查、风险预警、问题处置 \\
\bottomrule
\end{tabular}
\end{table}

\subsection{贷后管理的"三查"制度}

贷后管理是银行风控的最后一道防线，"三查"制度是其核心：

\begin{itemize}[leftmargin=2em]
  \item \textbf{贷前调查}（Investigation）：银行在授信前对借款人的经营状况、财务状况、还款能力、担保措施进行实地调查和核实，确保"了解你的客户"
  \item \textbf{贷时审查}（Review）：在放款时对合同条款、担保登记、放款条件进行逐项审核，确保"条件全落实"
  \item \textbf{贷后检查}（Monitoring）：贷款发放后定期跟踪借款人经营和财务变化，及时发现风险信号，确保"动态可控"
\end{itemize}

\begin{theorybox}[信息不对称理论与银行信贷]
乔治·阿克洛夫（George Akerlof）、迈克尔·斯彭斯（Michael Spence）和约瑟夫·斯蒂格利茨（Joseph Stiglitz）因信息不对称理论获得2001年诺贝尔经济学奖。该理论对银行信贷有深刻解释力：

\textbf{1. 逆向选择}（Adverse Selection）：在贷款审批前，借款人比银行更了解自身的真实风险状况。高风险借款人更愿意接受高利率，而低风险借款人可能因利率过高而退出市场，导致银行贷款组合整体风险上升。

\textbf{2. 道德风险}（Moral Hazard）：贷款发放后，借款人可能将资金用于比约定风险更高的项目——如果项目成功，借款人获得超额收益；如果失败，损失主要由银行承担。这种"收益私有化、风险社会化"的激励扭曲是银行信贷风险的核心来源。

\textbf{银行的应对机制}：
\begin{itemize}[leftmargin=1em]
  \item \textbf{抵质押担保}：提高借款人违约成本，抑制道德风险
  \item \textbf{限制性条款}（Covenant）：在贷款合同中约定财务指标约束
  \item \textbf{分期放款}：根据项目进度分批放款，降低一次性风险暴露
  \item \textbf{贷后监控}：定期获取财务报表和经营信息，缓解信息不对称
\end{itemize}

AI技术为缓解信息不对称提供了新工具：通过大数据分析企业关联关系、资金流向和行为异常，可以更早、更全面地识别风险信号。
\end{theorybox}

\subsection{中小企业融资难的理论分析}

中小企业融资难是全球性难题，其根源在于信息不对称在中小企业信贷中更为严重：

\begin{enumerate}[leftmargin=2em]
  \item \textbf{财务信息不透明}：中小企业财务报表规范性差，银行难以准确评估其还款能力
  \item \textbf{缺乏合格抵质押物}：中小企业资产以存货、应收账款为主，传统银行偏好不动产抵押
  \item \textbf{规模不经济}：中小企业贷款金额小、审批流程与大企业相同，单笔服务成本高
  \item \textbf{风险溢价过高}：银行为弥补信息不对称带来的风险，要求更高的利率，形成恶性循环
\end{enumerate}

AI和大数据技术正在改变这一困境：通过分析企业的税务数据、发票数据、水电费等"替代数据"（Alternative Data），银行可以在缺乏传统财务信息的情况下对中小企业进行信用评估，降低信息不对称程度，从而扩大信贷覆盖面。

% ============================================================
\section{企业信用风险评估模型}
% ============================================================

\subsection{Altman Z-Score模型}

Altman Z-Score模型由纽约大学Edward Altman教授于1968年提出，是最经典的财务困境预测模型。该模型通过多元判别分析（MDA），将5个财务比率加权求和，得到Z值用于判断企业的破产风险。

原始Z-Score公式（适用于上市公司）：

\[
Z = 1.2 X_1 + 1.4 X_2 + 3.3 X_3 + 0.6 X_4 + 1.0 X_5
\]

其中各变量的定义和经济含义如下：

\begin{table}[H]
\centering
\begin{tabular}{cllp{5cm}}
\toprule
\textbf{变量} & \textbf{计算公式} & \textbf{经济含义} \\
\midrule
$X_1$ & 营运资本/总资产 & 衡量企业短期流动性和资产利用效率 \\
$X_2$ & 留存收益/总资产 & 衡量企业累积盈利能力和经营年限 \\
$X_3$ & 息税前利润/总资产 & 衡量企业真实盈利能力（扣除融资和税影响前） \\
$X_4$ & 股东权益市值/总负债 & 衡量企业偿债能力和市场信心 \\
$X_5$ & 销售收入/总资产 & 衡量企业资产周转效率 \\
\bottomrule
\end{tabular}
\end{table}

Z值的判定区间：

\begin{table}[H]
\centering
\begin{tabular}{llp{6cm}}
\toprule
\textbf{Z值范围} & \textbf{风险等级} & \textbf{含义} \\
\midrule
$Z > 2.99$ & 安全区 & 企业财务状况健康，破产概率极低 \\
$1.81 < Z \leq 2.99$ & 灰色区 & 财务状况不确定，需进一步分析 \\
$Z \leq 1.81$ & 危险区 & 企业面临严重的财务困境风险 \\
\bottomrule
\end{tabular}
\end{table}

\textbf{中国制造业适配版本}：由于中国会计准则和市场环境与美国不同，学者们对Z-Score进行了本土化修正。中国制造业适用的Z值判定标准通常调整为：

\begin{itemize}[leftmargin=2em]
  \item $Z > 2.675$：安全区
  \item $1.23 < Z \leq 2.675$：灰色区
  \item $Z \leq 1.23$：危险区
\end{itemize}

对于非上市企业，$X_4$中的"股东权益市值"替换为"股东权益账面价值"，权重系数也需相应调整。Altman针对非上市企业提出了Z'模型：

\[
Z' = 0.717 X_1 + 0.847 X_2 + 3.107 X_3 + 0.420 X_4' + 0.998 X_5
\]

其中$X_4' = $股东权益账面价值/总负债。

\subsection{KMV模型简介}

KMV模型由KMV公司（现Moody's Analytics）开发，基于Merton（1974）的期权定价理论，将企业股权视为以企业资产为标的的看涨期权：

\begin{itemize}[leftmargin=2em]
  \item 当企业资产价值$>$负债面值时，股东"行权"偿还债务，保留剩余价值
  \item 当企业资产价值$<$负债面值时，股东"放弃行权"，企业违约
  \item 违约距离（Distance to Default, DD）=（资产价值 - 违约点）/（资产价值 $\times$ 资产波动率）
\end{itemize}

KMV模型的优势在于利用股票市场的实时信息，具有前瞻性；缺点是仅适用于上市公司，且对资产价值波动率的估计敏感。

\subsection{逻辑回归评分模型}

逻辑回归（Logistic Regression）是银行信用评分中最广泛使用的统计模型。其建模流程如下：

\begin{enumerate}[leftmargin=2em]
  \item \textbf{变量选择}：从候选变量中筛选出对违约有显著预测力的变量，常用方法包括逐步回归、LASSO等
  \item \textbf{WOE编码}：将每个变量的取值区间转化为WOE值（Weight of Evidence），公式为：
  \[
  \text{WOE}_i = \ln\left(\frac{\text{好客户占比}_i}{\text{坏客户占比}_i}\right)
  \]
  WOE值越大，说明该区间好客户比例越高
  \item \textbf{IV值筛选}：计算每个变量的信息价值（Information Value），用于评估变量的预测能力：
  \[
  \text{IV} = \sum_{i=1}^{n} (\text{好客户占比}_i - \text{坏客户占比}_i) \times \text{WOE}_i
  \]
  一般认为IV $>$ 0.3的变量具有较强预测力
  \item \textbf{逻辑回归建模}：以WOE化后的变量为自变量，违约/正常为因变量，建立逻辑回归模型
  \item \textbf{模型验证}：使用KS统计量、AUC-ROC曲线等指标评估模型区分度
\end{enumerate}

\begin{theorybox}[三种信用风险模型的适用场景对比]
\begin{tabular}[t]{lp{3.5cm}p{3.5cm}p{3.5cm}}
  & \textbf{Altman Z-Score} & \textbf{KMV模型} & \textbf{逻辑回归评分卡} \\
\midrule
数据需求 & 财务报表 & 股价+财务数据 & 财务+行为+征信数据 \\
适用对象 & 上市/非上市企业 & 上市公司 & 企业+个人 \\
时效性 & 滞后（季度报表） & 实时（日频股价） & 中等（月度更新） \\
可解释性 & 强 & 中 & 强 \\
预测精度 & 中 & 较高 & 较高 \\
行业适用性 & 制造业最优 & 不限行业 & 需按行业建模 \\
\end{tabular}

\vspace{0.5em}
在银行实践中，通常组合使用多种模型：Z-Score用于快速筛查和持续监控，KMV模型用于上市公司组合管理，逻辑回归评分卡用于标准化授信审批。
\end{theorybox}

\subsection{实操：用AI计算Altman Z-Score}

\begin{practicebox}[实操：Altman Z-Score信用风险评估]
\textbf{工具}：finance-expert Skill + AI对话

\textbf{完整提示词}：

\begin{lstlisting}
请使用 finance-expert 技能完成企业信用风险评估：

【企业基本信息】
企业名称：恒达机械制造有限公司
行业：通用设备制造业
成立时间：2012年
员工人数：156人

【财务数据（2025年度）】
单位：万元
- 流动资产：1,800
- 流动负债：1,200
- 总资产：2,800
- 总负债：1,600
- 股东权益（账面价值）：1,200
- 留存收益：480
- 息税前利润（EBIT）：310
- 销售收入：4,100
- 净利润：250

【任务1：计算原始Z-Score】
请按以下公式计算：
Z = 1.2*(X1) + 1.4*(X2) + 3.3*(X3) + 0.6*(X4) + 1.0*(X5)

其中：
X1 = 营运资本/总资产 = (流动资产-流动负债)/总资产
X2 = 留存收益/总资产
X3 = 息税前利润/总资产
X4 = 股东权益账面价值/总负债
X5 = 销售收入/总资产

展示每一步计算过程。

【任务2：判定风险等级】
基于中国制造业适配标准：
- Z > 2.675 -> 安全区
- 1.23 <= Z <= 2.675 -> 灰色区
- Z < 1.23 -> 危险区

【任务3：敏感性分析】
如果销售收入下降20%（从4100万降至3280万），Z值如何变化？
如果总负债增加30%（从1600万增至2080万），Z值如何变化？
两种情况同时发生呢？

【任务4：生成评估报告】
输出结构化报告，包含：
1. 企业画像概要（行业/规模/成立年限）
2. Z-Score计算明细（五项指标逐项展示）
3. 风险等级判定及依据
4. 敏感性分析结果
5. 风险缓释建议（至少3条具体措施）

保存到 reports/altman_zscore_hengda.md
\end{lstlisting}

\textbf{学生练习}：修改财务数据，观察Z值和风险等级的变化。尝试以下场景：
\begin{itemize}[leftmargin=1em]
  \item 场景1：应收账款周转困难，流动资产降至1200万元
  \item 场景2：企业盈利改善，EBIT从310万增至500万
  \item 场景3：企业举债扩张，总负债从1600万增至2400万
\end{itemize}
\end{practicebox}

% ============================================================
\section{贷款定价模型}
% ============================================================

\subsection{FTP转移定价理论}

内部资金转移定价（Funds Transfer Pricing, FTP）是银行内部资金管理中心与各业务条线之间的资金价格。它的核心作用是：

\begin{enumerate}[leftmargin=2em]
  \item \textbf{分离利率风险}：将利率风险从业务条线集中到资金中心统一管理
  \item \textbf{公平绩效考核}：剔除资金成本差异后，真实反映各业务条线的盈利贡献
  \item \textbf{引导定价决策}：为贷款定价和存款定价提供成本基准
\end{enumerate}

\subsection{贷款定价公式}

银行贷款定价的基本公式为：

\[
\text{贷款利率} = \text{FTP} + \text{风险溢价} + \text{运营成本} + \text{资本占用成本} + \text{目标利润}
\]

各组成部分的经济含义：

\begin{table}[H]
\centering
\begin{tabular}{lp{9cm}}
\toprule
\textbf{定价要素} & \textbf{含义与计算方法} \\
\midrule
FTP（资金成本） & 银行内部资金转移价格，与贷款期限匹配。由资金中心根据市场利率定期公布 \\
风险溢价 & 基于客户信用等级的差异化定价，反映预期违约损失（EL = PD $\times$ LGD $\times$ EAD） \\
运营成本 & 人力成本、系统成本、网点分摊等，通常按贷款金额的固定比例计收 \\
资本占用成本 & 贷款占用监管资本的代价 = 贷款金额 $\times$ 风险权重 $\times$ 资本充足率 $\times$ 资本回报率 \\
目标利润 & 银行股东要求的最低回报率 \\
\bottomrule
\end{tabular}
\end{table}

\subsection{风险溢价的确定}

风险溢价是贷款定价中最核心也最难以量化的部分，其本质是对预期信用损失的补偿：

\begin{table}[H]
\centering
\begin{tabular}{lcc}
\toprule
\textbf{信用等级} & \textbf{年化违约概率（PD）} & \textbf{建议风险溢价} \\
\midrule
AAA & 0.05\% & 0.30\% \\
AA & 0.20\% & 0.60\% \\
A & 0.50\% & 1.00\% \\
BBB & 1.50\% & 1.80\% \\
BB & 4.00\% & 3.00\% \\
B & 8.00\% & 5.00\% \\
\bottomrule
\end{tabular}
\end{table}

\subsection{资本占用成本与巴塞尔协议III}

根据巴塞尔协议III的标准法，不同类别贷款的风险权重不同：

\begin{itemize}[leftmargin=2em]
  \item 对公贷款（一般企业）：风险权重100\%
  \item 对公贷款（中小企业）：风险权重85\%
  \item 个人住房抵押贷款：风险权重50\%
  \item 个人消费贷款：风险权重75\%
  \item 地方政府债券：风险权重20\%--50\%
\end{itemize}

资本占用成本的计算公式：

\[
\text{资本占用成本} = \text{贷款金额} \times \text{风险权重} \times \text{资本充足率} \times \text{资本回报率}
\]

例如：一笔500万元的一般企业贷款，风险权重100\%，资本充足率12.5\%，资本回报率12\%：

\[
\text{资本占用成本率} = 100\% \times 12.5\% \times 12\% = 1.50\%
\]

\subsection{实操：用xlsx Skill构建贷款定价Excel模型}

\begin{practicebox}[实操：银行贷款FTP定价模型]
\textbf{工具}：xlsx Skill

\textbf{完整提示词}：

\begin{lstlisting}
请使用 xlsx 技能创建银行贷款定价模型 Excel 文件：

【文件名】reports/loan_pricing_model.xlsx

【Sheet 1：定价参数】
构建贷款定价公式：
贷款利率 = 资金成本(FTP) + 风险溢价 + 运营成本 + 资本占用成本 + 目标利润

参数设置（使用 Excel 公式，不硬编码计算结果）：
- 资金成本（FTP）：一年期 2.8%，三年期 3.2%
- 风险溢价：按信用等级差异化
  * AAA级：0.3%  AA级：0.6%  A级：1.0%  BBB级：1.8%  BB级：3.0%
- 运营成本：0.4%（含人工、系统、网点分摊）
- 资本占用成本：贷款金额 * 风险权重(100%) * 资本充足率(12.5%) * 资本回报率(12%)
- 目标利润：0.2%

【Sheet 2：客户定价测算】
为 5 位不同客户自动计算贷款报价：
| 客户名 | 金额(万) | 期限 | 信用等级 | 抵押方式 | 最终利率 |
| 恒达机械 | 500 | 1年 | A | 厂房抵押 | =公式计算 |
| 绿源农业 | 200 | 3年 | BBB | 信用 | =公式计算 |
| 华新科技 | 1000 | 1年 | AA | 应收质押 | =公式计算 |
| 鑫达贸易 | 300 | 1年 | BB | 保证担保 | =公式计算 |
| 瑞丰地产 | 2000 | 3年 | A | 土地抵押 | =公式计算 |

要求：
- FTP根据期限用VLOOKUP从Sheet1查找
- 风险溢价根据信用等级用VLOOKUP从Sheet1查找
- 资本占用成本用公式计算
- 最终利率 = 各项之和

【Sheet 3：敏感性分析】
- 当 FTP 上浮/下浮 20bp 时，各客户利率如何变化
- 当信用等级上调/下调一级时，利率如何变化
- 用条件格式标注：利率 < 4% 绿色，4-6% 黄色，> 6% 红色

【Sheet 4：定价决策表】
生成综合决策表：
- 每位客户的基准利率、最优利率（信用等级上调一级）、最差利率（信用等级下调一级）
- 利率弹性：FTP每变动10bp，利率变动多少bp
- 利润测算：每笔贷款的年化利润额

【格式要求】
- 所有计算必须使用 Excel 公式（=SUM, =VLOOKUP, =IF 等）
- 表头加粗 + 蓝色背景
- 百分比统一保留 2 位小数
- 金额使用千位分隔符
\end{lstlisting}
\end{practicebox}

% ============================================================
\section{贷后风险预警系统}
% ============================================================

\subsection{早期预警理论}

贷后风险预警（Early Warning System, EWS）是银行在贷款存续期内，通过监测一系列先行指标，在客户违约前及时发现风险信号并采取缓释措施的管理系统。

经济指标按与经济周期的关系可分为三类：

\begin{itemize}[leftmargin=2em]
  \item \textbf{先行指标}（Leading Indicators）：在风险事件发生前出现变化，如订单下滑、库存增加、行业政策收紧等——这是预警系统的核心监测对象
  \item \textbf{同步指标}（Coincident Indicators）：与风险事件同时变化，如营收下降、现金流恶化等
  \item \textbf{滞后指标}（Lagging Indicators）：在风险事件发生后才显现，如逾期、不良贷款——此时损失已产生，预警价值有限
\end{itemize}

有效的预警系统应聚焦先行指标和同步指标，在损失发生前发出信号。

\subsection{预警信号分类}

\begin{table}[H]
\centering
\begin{tabular}{lp{5cm}p{5cm}}
\toprule
\textbf{信号类别} & \textbf{典型信号} & \textbf{监测方式} \\
\midrule
财务预警 & 现金流持续为负、应收账款激增$>$30\%、存货积压、资产负债率突破70\% & 季度报表分析+AI自动比对 \\
行为预警 & 贷款资金挪用、关联交易异常、实控人变更、大额资金转出 & 资金流水监控+AI异常检测 \\
外部预警 & 行业下行周期、环保处罚、诉讼信息、税务异常、高管被调查 & 舆情监控+工商/司法数据接口 \\
\bottomrule
\end{tabular}
\end{table}

\subsection{三色预警体系}

银行通常采用三色预警体系对风险信号进行分级管理：

\begin{table}[H]
\centering
\begin{tabular}{lp{3cm}p{4cm}p{4cm}}
\toprule
\textbf{预警级别} & \textbf{触发条件} & \textbf{处置要求} & \textbf{上报时限} \\
\midrule
黄色预警（关注） & 1--2个轻度预警信号 & 客户经理7日内实地走访核实 & 7日内 \\
橙色预警（警告） & 2个以上轻度或1个中度信号 & 立即电话联系，3日内上报风控部门 & 3日内 \\
红色预警（紧急） & 重大风险信号（逾期/实控人失联/诉讼） & 启动应急预案，当日上报分管行长 & 当日 \\
\bottomrule
\end{tabular}
\end{table}

\subsection{实操：用bank-risk-alert Skill生成预警通知}

\begin{practicebox}[实操：贷后风险预警话术生成]
\textbf{工具}：bank-risk-alert Skill（第4章实验中自写的Skill）

如果尚未安装该Skill，请先创建：

\begin{lstlisting}[style=shell]
# 在 skills 目录下创建 bank-risk-alert Skill
mkdir -p skills/bank-risk-alert
\end{lstlisting}

\textbf{预警测试提示词}（完整可用）：

\begin{lstlisting}
请使用 bank-risk-alert 技能生成风险预警通知：

【预警场景1：橙色预警】
- 客户名称：恒达机械制造有限公司
- 预警信号：连续2期利息逾期 + 应收账款周转天数从60天延长至120天
- 风险等级：橙色预警
- 贷款余额：500万元
- 贷款到期日：2026-09-30
- 客户经理：张经理

请生成以下内容：
1. 预警摘要卡片（一屏可看完的关键信息）
2. 客户经理电话话术（开场白->风险告知->解决方案->行动承诺->结束语）
3. 内部上报邮件模板（收件人=分管行长）
4. 跟进工单（任务项+负责人+截止日期+检查标准）

【预警场景2：红色预警】
- 客户名称：绿源农业科技有限公司
- 预警信号：实控人失联 + 公司账户资金异常转出300万元 + 税务部门处罚通知
- 风险等级：红色预警
- 贷款余额：200万元
- 贷款到期日：2026-12-31
- 客户经理：李经理

请生成完整的红色预警响应方案，包含：
1. 紧急预警通知
2. 应急处置方案（冻结账户/保全资产/法务介入）
3. 客户经理上门话术
4. 内部上报邮件
5. 后续处置时间表
\end{lstlisting}

\textbf{电话话术模板参考}（橙色预警场景）：

\begin{lstlisting}
请为以下贷后预警场景生成客户经理电话话术：

【客户】恒达机械制造有限公司 王总
【预警】连续2期利息逾期（每月应还利息约2.3万元）
【目标】了解逾期原因 + 督促还款 + 评估是否需要追加担保

话术结构：
1. 开场白（自我介绍+关心询问，不直接说"你逾期了"）
2. 风险告知（委婉提醒有两期利息未按时支付）
3. 原因了解（了解企业经营是否出现困难）
4. 解决方案（提供还款安排/展期/追加担保等选项）
5. 行动承诺（确认具体还款日期和金额）
6. 结束语（表达支持+明确后续联系时间）

语气要求：专业但不咄咄逼人，体现银行对客户的关心
\end{lstlisting}

\textbf{内部上报邮件模板参考}：

\begin{lstlisting}
请生成贷后预警内部上报邮件：

收件人：分管行长
抄送：风险管理部总经理、支行行长
主题：【橙色预警】恒达机械制造有限公司风险预警报告

邮件内容包含：
1. 企业基本信息（名称/贷款余额/到期日/担保方式）
2. 预警信号描述（具体事实+数据）
3. 已采取措施（客户经理已做的工作）
4. 风险评估（当前风险等级+发展趋势判断）
5. 处置建议（追加担保/降额/提前收回/展期等）
6. 时间要求（请在XX日前批复）
\end{lstlisting}
\end{practicebox}

% ============================================================
\section{行业风险与供应链金融分析}
% ============================================================

\subsection{行业风险分析框架}

银行在进行对公授信时，不仅要评估单个企业的风险，还必须考量其所在行业的系统性风险。行业风险分析框架包含以下维度：

\begin{table}[H]
\centering
\begin{tabular}{lp{5cm}p{5cm}}
\toprule
\textbf{分析维度} & \textbf{评估内容} & \textbf{风险信号} \\
\midrule
行业生命周期 & 导入期/成长期/成熟期/衰退期 & 衰退期行业整体信贷紧缩 \\
行业集中度 & CR3/CR5/HHI指数 & 过度集中=系统性风险 \\
政策敏感性 & 监管政策/产业政策/环保政策 & 政策收紧行业面临压缩退出 \\
周期性 & 与GDP周期的相关度 & 强周期行业违约率波动大 \\
技术替代性 & 新技术颠覆的可能性 & 技术替代导致资产价值归零 \\
\bottomrule
\end{tabular}
\end{table}

\subsection{供应链金融模式}

供应链金融（Supply Chain Finance）是银行围绕核心企业，管理上下游中小企业的资金流和物流，把单个企业的不可控风险转变为供应链企业整体的可控风险。三种经典模式：

\begin{table}[H]
\centering
\begin{tabular}{lp{4cm}p{3cm}p{4cm}}
\toprule
\textbf{模式} & \textbf{运作方式} & \textbf{核心风控} & \textbf{适用场景} \\
\midrule
应收账款融资 & 供应商将对核心企业的应收账款转让给银行融资 & 核心企业确权 & 上游供应商融资 \\
预付款融资 & 经销商预付货款后，银行控制提货权 & 货物监管+回购承诺 & 下游经销商融资 \\
存货融资 & 企业以存货作为质押物获取贷款 & 第三方仓储监管 & 有大量存货的企业 \\
\bottomrule
\end{tabular}
\end{table}

\subsection{AI在供应链金融中的应用}

AI技术正在从三个层面重塑供应链金融：

\begin{enumerate}[leftmargin=2em]
  \item \textbf{发票验真}：利用OCR+NLP自动识别和验证发票信息，比对税务系统数据，防止虚假发票融资诈骗。传统人工验真单张发票需3--5分钟，AI系统可在1秒内完成
  \item \textbf{物流监控}：通过物联网传感器+GPS实时监控质押物状态（位置、温度、数量），确保质押物真实存在且未被挪用。一旦出现异常移动，系统自动触发预警
  \item \textbf{动态额度}：基于企业的实时经营数据（订单、发货、回款），AI动态调整授信额度，替代传统的静态授信模式。企业经营好则额度自动上调，经营恶化则自动下调，实现"用多少贷多少"
\end{enumerate}

\begin{theorybox}[核心企业信用传导机制]
供应链金融的核心逻辑是"信用传导"：核心企业（通常是行业龙头）的优质信用，通过贸易关系"传导"给上下游的中小企业，使其能够以更低的成本获得融资。

信用传导的路径：
\begin{enumerate}[leftmargin=1em]
  \item 核心企业向供应商采购，产生应付账款
  \item 供应商持有核心企业的应付账款（应收账款），本质上是持有核心企业的信用
  \item 银行基于核心企业的信用，为供应商提供应收账款融资
  \item 核心企业的信用从"间接"变为"直接"——银行实质上是在对核心企业信用进行定价
\end{enumerate}

信用传导的关键风险在于"传导断裂"：如果核心企业否认债务、延迟付款或自身出现风险，信用传导链条就会断裂，导致整个供应链融资体系崩塌。2021年某大型房地产企业违约事件，就是信用传导断裂的典型案例——其上下游数百家中小企业同时面临融资困难。
\end{theorybox}

% ============================================================
\section{实验：中小企业500万贷款全流程审批模拟}
% ============================================================

本实验将串联本章所有知识点，模拟一笔中小企业500万元流动资金贷款从申请到贷后管理的完整流程。

\begin{exercisebox}[实验：中小企业500万贷款全流程审批模拟]

\textbf{实验目标}：体验一笔对公贷款从申请到贷后管理的全流程，掌握信用评估、定价、审批、预警四大核心技能

\textbf{实验时间}：120分钟

\textbf{企业背景}：

恒达机械制造有限公司，成立于2012年，主营通用机械零部件加工，员工156人，位于四川省成都市某工业园区。现向某城商行申请500万元流动资金贷款，期限12个月，用途为原材料采购。

\textbf{Step 1：企业基本信息收集}（15分钟）

\begin{lstlisting}
请使用 finance-expert 技能完成企业基本信息收集：

【任务】为恒达机械制造有限公司建立信贷档案

1. 行业分析：
   - 通用设备制造业的行业生命周期阶段
   - 行业集中度与竞争格局
   - 近3年行业景气度趋势
   - 主要政策风险

2. 企业画像：
   - 经营年限：14年（中等偏上）
   - 规模：小型企业（营收4100万）
   - 员工：156人
   - 主营业务：机械零部件加工
   - 主要客户：3家大型制造企业
   - 供应商：原材料采购依赖2家主要供应商

3. 初步风险评估：
   - 行业风险等级
   - 经营稳定性评估
   - 关联关系排查（是否存在集团关联）

输出为结构化的企业调查摘要
\end{lstlisting}

\textbf{Step 2：财务报表分析与Z-Score计算}（25分钟）

\begin{lstlisting}
请使用 finance-expert 技能完成财务分析与Z-Score计算：

【企业财务数据（2023-2025年）】
单位：万元

利润表：
           2023    2024    2025
营业收入   3,200   3,800   4,100
营业成本   2,560   2,960   3,200
EBIT       250     310     310
净利润     180     220     250

资产负债表（2025年末）：
流动资产   1,800
其中：货币资金    300
      应收账款    800
      存货        600
流动负债   1,200
总资产     2,800
总负债     1,600
股东权益   1,200
留存收益   480

【分析要求】
1. 三年财务趋势分析（营收/利润/资产负债率变化）
2. 关键财务指标计算：
   - 资产负债率 = 总负债/总资产
   - 流动比率 = 流动资产/流动负债
   - 速动比率 = (流动资产-存货)/流动负债
   - 应收账款周转天数 = 应收账款/(营业收入/365)
   - 存货周转天数 = 存货/(营业成本/365)
   - 利息保障倍数 = EBIT/利息支出（假设利率5.5%）
3. Altman Z'-Score计算（非上市企业版本）
4. DuPont分析：ROE分解
5. 贷款偿还能力评估：经营性现金流/贷款本息

输出完整的财务分析报告
\end{lstlisting}

\textbf{Step 3：FTP定价模型计算报价}（25分钟）

\begin{lstlisting}
请使用 xlsx 技能完成贷款定价计算：

【定价参数】
- 客户：恒达机械制造有限公司
- 贷款金额：500万元
- 贷款期限：12个月
- 信用等级：A级（基于Step2的Z-Score和财务分析）
- 抵押方式：厂房抵押（评估价值600万元，抵押率60%）
- FTP：一年期2.8%
- 风险溢价：A级=1.0%
- 运营成本：0.4%
- 资本占用：500万*100%*12.5%*12%=1.50%
- 目标利润：0.2%

【计算要求】
1. 计算贷款利率
2. 计算年利息收入
3. 计算抵押物可覆盖金额（600万*60%）
4. 计算风险敞口（贷款金额-抵押物覆盖金额）
5. 如果信用等级为BBB级，利率和利润如何变化？
6. 如果贷款金额降至300万（抵押充足覆盖），利率和利润如何变化？

保存到 reports/loan_pricing_hengda.xlsx
\end{lstlisting}

\textbf{Step 4：风险等级判定与审批决策}（20分钟）

\begin{lstlisting}
基于以上分析结果，请完成审批决策：

【综合评估】
1. 行业风险：通用设备制造业，成熟期，竞争充分，政策风险中等
2. 经营风险：经营14年，客户集中度偏高，供应商依赖度偏高
3. 财务风险：资产负债率57%，Z'-Score在灰色区，现金流偏紧
4. 担保风险：厂房抵押，评估价值600万，抵押率60%，覆盖360万

【审批决策模板】
请输出完整的审批意见书，包含：
1. 申请信息摘要
2. 风险评估结论（综合风险等级：低/中/高）
3. 审批决定（批准/有条件批准/拒绝）
4. 贷款条件（如需追加担保/限制性条款等）
5. 贷后管理要求（检查频次/重点监控指标）
6. 敏感性分析（什么情况下需要重新评估）
\end{lstlisting}

\textbf{Step 5：贷后预警方案设计}（20分钟）

\begin{lstlisting}
请使用 bank-risk-alert 技能设计贷后预警方案：

【贷款信息】
- 客户：恒达机械制造有限公司
- 贷款金额：500万元
- 到期日：2027-06-08
- 担保方式：厂房抵押

【预警方案要求】
1. 设定监测指标及阈值：
   - 财务指标：资产负债率>70%触发黄色预警，>80%触发橙色预警
   - 行为指标：利息逾期立即触发橙色预警，本金逾期触发红色预警
   - 外部指标：环保处罚/重大诉讼触发橙色预警

2. 检查频次：
   - 正常类：每季度实地检查一次
   - 关注类：每月电话+每季度实地
   - 不良类：每周跟踪

3. 生成一份标准化的季度贷后检查报告模板
\end{lstlisting}

\textbf{Step 6：生成审批意见书}（15分钟）

\begin{lstlisting}
请使用 docx 技能生成正式的贷款审批意见书：

【文件名】reports/loan_approval_hengda.docx

【文档结构】
1. 封面：XX银行授信审批意见书 + 日期 + 机密
2. 申请信息：
   - 借款人：恒达机械制造有限公司
   - 申请额度：500万元
   - 期限：12个月
   - 用途：原材料采购流动资金
3. 企业概况（Step1的摘要）
4. 财务分析（Step2的关键结论）
5. 定价方案（Step3的计算结果）
6. 风险评估（Step4的审批决定）
7. 贷后管理要求（Step5的预警方案）
8. 审批签章页（审批人签字栏）

【格式要求】
- A4竖版，银行商务蓝色主题
- 页眉：XX银行授信审批意见书（内部机密）
- 页脚：第X页/共Y页
- 标题层级清晰，使用正式公文语体
\end{lstlisting}

\textbf{预期产出清单}：
\begin{enumerate}[leftmargin=2em]
  \item 企业调查摘要（Markdown格式）
  \item 财务分析报告（含Z-Score计算过程）
  \item \texttt{reports/loan\_pricing\_hengda.xlsx}：定价计算模型
  \item 审批意见书（Markdown格式）
  \item 贷后预警方案（含检查报告模板）
  \item \texttt{reports/loan\_approval\_hengda.docx}：正式审批意见书
\end{enumerate}

\end{exercisebox}

% ============================================================
\section{案例研究：城商行数字化风控转型}
% ============================================================

\begin{casebox}[某城商行的风控升级之路]

\textbf{一、传统人工审批的痛点}

某中部地区城商行（资产规模约3000亿元），在数字化转型前面临以下风控困境：

\begin{itemize}[leftmargin=2em]
  \item \textbf{审批效率低}：一笔对公贷款从受理到审批平均需要25个工作日，客户等待时间长，部分优质客户因此流失
  \item \textbf{风控标准不统一}：不同支行、不同审批人的风险偏好差异大，同类客户可能得到截然不同的审批结果
  \item \textbf{贷后管理薄弱}：客户经理每人管理50--80户对公客户，难以实现有效的贷后检查，往往是"出了问题才发现"
  \item \textbf{数据利用率低}：行内积累了大量客户数据，但散落在不同系统（信贷系统、核心系统、CRM），无法有效整合分析
\end{itemize}

\textbf{二、AI风控系统上线过程}

该行于2023年启动"智慧风控"项目，分三阶段推进：

\begin{enumerate}[leftmargin=2em]
  \item \textbf{数据治理阶段}（6个月）：整合信贷、核心、外管、征信四大系统数据，建立统一客户画像数据库，覆盖对公客户1.2万户
  \item \textbf{模型建设阶段}（8个月）：与金融科技公司合作，构建信用评分模型、反欺诈模型和预警模型，基于5年历史数据训练验证
  \item \textbf{系统部署阶段}（4个月）：将AI模型嵌入信贷审批流程，实现"人机协同"审批——AI给出评分和建议，审批人做最终决策
\end{enumerate}

\textbf{三、效果数据}

系统上线一年后，关键指标显著改善：

\begin{table}[H]
\centering
\begin{tabular}{lcc}
\toprule
\textbf{指标} & \textbf{上线前} & \textbf{上线后} \\
\midrule
对公贷款平均审批时间 & 25个工作日 & 8个工作日 \\
审批标准一致性 & 低（主观判断为主） & 高（AI评分+人工复核） \\
贷后预警及时率 & 约30\% & 约85\% \\
新生成不良贷款率 & 2.1\% & 1.4\% \\
客户经理人均管户数 & 60户 & 85户 \\
\end{tabular}
\end{table}

\textbf{四、对金融专业学生的启示}

\begin{enumerate}[leftmargin=2em]
  \item \textbf{AI不会取代风控人员，但会用AI的风控人员会取代不会的}：系统上线后，风控人员的工作从"手工审批"转向"模型监控+异常处理"，对复合型人才的需求更加强烈
  \item \textbf{数据质量是AI风控的基石}：该行花了6个月做数据治理，占总项目时间的1/3——"Garbage in, garbage out"对AI尤为致命
  \item \textbf{人机协同是当前最佳实践}：完全依赖AI审批存在合规风险，完全依赖人工又无法提升效率，两者结合是银行的现实选择
  \item \textbf{金融专业知识是AI落地的关键}：AI模型的设计、验证和监控都需要深厚的金融专业背景，纯技术人员无法独立完成
\end{enumerate}
\end{casebox}

% ============================================================
% 本章小结
% ============================================================

\section*{本章小结}

本章系统介绍了AI在银行公司业务和风控建模中的应用，核心内容包括：

\begin{enumerate}[leftmargin=2em]
  \item \textbf{公司业务理论}：公司银行的核心是授信全流程管理，信息不对称理论（逆向选择+道德风险）是理解银行信贷风险的基石
  \item \textbf{信用风险评估}：Altman Z-Score、KMV模型和逻辑回归评分卡各有适用场景，银行通常组合使用。AI技术使信用评估更高效，但可解释性仍是监管关注的重点
  \item \textbf{贷款定价}：FTP定价公式（贷款利率 = FTP + 风险溢价 + 运营成本 + 资本占用 + 目标利润）是银行贷款定价的基本框架，xlsx Skill可以高效构建定价模型
  \item \textbf{贷后预警}：三色预警体系（黄/橙/红）是贷后风险管理的核心机制，bank-risk-alert Skill可以自动生成预警话术和通知
  \item \textbf{供应链金融}：核心企业信用传导是供应链金融的核心逻辑，AI在发票验真、物流监控和动态额度方面展现了巨大价值
\end{enumerate}

\keyterms{
公司银行、授信全流程、三查制度、信息不对称、逆向选择、道德风险、Altman Z-Score、KMV模型、逻辑回归、WOE、IV值、FTP定价、风险溢价、资本占用、巴塞尔协议III、贷后预警、三色预警、供应链金融、信用传导
}

\begin{exercisebox}[思考题]
\begin{enumerate}
  \item 信息不对称理论如何解释中小企业融资难问题？AI大数据技术能在多大程度上缓解信息不对称？
  \item 某企业Z-Score为2.1（处于灰色区），但其行业处于衰退期且应收账款持续增加。你会如何综合判断其信用风险？是否应该单纯依赖Z-Score的结果？
  \item 在贷款定价公式中，为什么资本占用成本要计入贷款利率？如果银行资本充足率从12.5\%提高到15\%，贷款利率会如何变化？
  \item 为什么供应链金融中核心企业的信用可以"传导"给上下游中小企业？这种信用传导在什么条件下可能失效？
  \item \textbf{案例分析题}：某城商行在对一家制造业企业审批500万元贷款时，AI风控系统给出的评分是"中风险"，但客户经理认为该企业"经营稳健、还款意愿强"，主张审批通过。你作为风险管理部审查员，如何处理这种人机判断不一致的情况？请提出具体的决策框架。
\end{enumerate}
\end{exercisebox}
